\subsection{与预层相伴的平展空间}

设 B 为一个拓扑空间，\(\mathcal{B}\) 为 B 的拓扑的一个基，\(\mathcal{F} = (\mathcal{F}(U), r_{UV})\) 为 B 上的、相对于基 \(\mathcal{B}\) 的预层。设 L 为对 \((U, s)\) （满足 \(U \in \mathcal{B}\) 且 \(s \in \mathcal{F}(U)\) 者）所成的集。用 \(X_{\mathcal{F}}\) 表示族 \((U)_{(U,s) \in L}\) 的和空间。于是 \(X_{\mathcal{F}}\) 是三元组 \((U, s, x)\) 所成的集，其中 \(U \in \mathcal{B}\) 、\(s \in \mathcal{F}(U)\) 、\(x \in U\) 。设 \(R_{\mathcal{F}}\) 为 \(X_{\mathcal{F}}\) 上的、由下式定义的关系：\(R_{\mathcal{F}}((U,s,x),(U',s',x'))\) 当且仅当“ \(x = x'\) 且存在 \(W \in \mathcal{B}\) 使 \(x \in W\) 、\(W \subset U \cap U'\) 且 \(r_{WU}(s) = r_{WU'}(s')\) ”。关系 \(R_{\mathcal{F}}\) 是 \(X_{\mathcal{F}}\) 中的一个等价关系：按定义它是自反的且对称的；我们来证明它是传递的。设 \(\xi = (U, s, x)\) 、\(\xi' = (U', s', x')\) 与 \(\xi'' = (U'', s'', x'')\) 为 \(X_{\mathcal{F}}\) 中满足 \(R_{\mathcal{F}}(\xi,\xi')\) 与 \(R_{\mathcal{F}}(\xi',\xi'')\) 的元素。于是有 \(x = x' = x''\) ，且存在两个元素 \(W'\) 与 \(W''\) （\(\mathcal{B}\) 中的，包含 \(x\) 者），使得 \(W' \subset U \cap U'\) 、\(W'' \subset U' \cap U''\) 、\(r_{W'U}(s) = r_{W'U'}(s')\) 、\(r_{W''U'}(s') = r_{W''U''}(s'')\) 。设 W 为 \(\mathcal{B}\) 中包含 \(x\) 且含于 \(W' \cap W''\) 的一个元素。于是有 \(W \subset U \cap U''\) ，

\[
r_{\mathrm{WU}}(s) = r_{\mathrm{WW'}} \circ r_{\mathrm{W'U}}(s) = r_{\mathrm{WW'}} \circ r_{\mathrm{W'U'}}(s') = r_{\mathrm{WU'}}(s')
\]

并且类似地有 \(r_{\mathrm{WU}'}(s') = r_{\mathrm{WU}''}(s'')\) 。因此有 \(R_{\mathcal{F}}(\xi,\xi'')\) ，从而关系 \(R_{\mathcal{F}}\) 是传递的。

用 \(E_{F}\) 表示商集 \(X_{F}/R_{F}\) ，并用 \([U,s,x]\) 表示 \(E_{F}\) 中的、元素 \((U,s,x)\) （\(X_{F}\) 的）的典范像。对 \(U \in B\) 与 \(s \in \mathcal{F}(U)\) ，用 \(\sigma_{\mathcal{F}}(U,s)\colon U \to E_{\mathcal{F}}\) 表示映射 \(x \mapsto [U,s,x]\) 。给集合 \(E_{F}\) 赋以商拓扑，即使得诸映射 \(\sigma_{\mathcal{F}}(U,s)\) （对 \(U \in B\) 与 \(s \in \mathcal{F}(U)\) ）都连续的最细拓扑。映射 \(pr_{3}\colon X_{F} \to B\) 过渡到商，定义了一个连续映射 \(p\colon E_{F} \to B\) ：有 \(p([U,s,x]) = x\) 。

命题 1. --- 映射 \(p\colon E_{\mathcal{F}} \to B\) 是平展的。于是，对每个开集 \(U \in \mathcal{B}\) 与每个 \(s \in \mathcal{F}(U)\) ，映射 \(\sigma_{\mathcal{F}}(U, s)\) 是 \(p\) 在 \(U\) 上的一个连续截面。

设 \(\lambda = (\mathrm{U}, s)\) 与 \(\mu = (\mathrm{U}', s')\) 为 L 的元素。按关系 \(\mathrm{R}_{\mathcal{F}}\) 的定义，集 \(\mathrm{A}_{\lambda \mu}\) （点 \(x\) （\(\mathrm{U} \cap \mathrm{U}'\) 的）中使 \(\sigma_{\mathcal{F}}(\mathrm{U}, s)\) 与 \(\sigma_{\mathcal{F}}(\mathrm{U}', s')\) 重合者所成）是由 \(x \in \mathrm{U} \cap \mathrm{U}'\) （满足 \(s(x) = s'(x)\) 者）构成的集在 B 中的内部。由此得出 \(\mathrm{A}_{\mu \lambda} = \mathrm{A}_{\lambda \mu}\) 。然后用 \(h_{\mu \lambda} \colon \mathrm{A}_{\lambda \mu} \to \mathrm{A}_{\mu \lambda}\) 表示映射 \(\mathrm{Id}_{\mathrm{A}_{\lambda \mu}}\) 。集合 \(\mathrm{E}_{\mathcal{F}}\) 是沿诸 \(\mathrm{A}_{\lambda \mu}\) 借助双射 \(h_{\mu \lambda}\) 粘合诸开集 U 而得的 (TG, I, p. 16)。根据 TG, I, p. 17 的 prop. 9，映射 \(\sigma_{\mathcal{F}}(\mathrm{U},s)\) 诱导从 U 到 \(\mathrm{E}_{\mathcal{F}}\) 的一个开子集上的一个同胚。这就证明了映射 \(p\) 是平展的 (I, p. 33, prop. 9)。

对每个开集 \(U \in B\) 与每个 \(s \in \mathcal{F}(U)\) ，有 \(\sigma_{\mathcal{F}}(U, s)(x) = [U, s, x]\) 对一切 \(x \in U\) 成立。于是第二个断言由 p 的定义得出。

定义 4. --- 上面定义的平展 B-空间 \((E_{\mathcal{F}},p)\) 称为与预层 \(\mathcal{F}\) 相伴的平展 B-空间。对 \(x \in \mathrm{B}\) ，\(E_{\mathcal{F}}\) 在 \(x\) 上的纤维称为预层 \(\mathcal{F}\) 在 \(x\) 处的茎，记作 \(\mathcal{F}_x\) 。对每个开集 \(U \in \mathcal{B}\) 、每个截面 \(s \in \mathcal{F}(\mathrm{U})\) （\(\mathcal{F}\) 在 U 上的）和 U 的每个点 \(x\) ，元素 \([U,s,x]\) （\(E_{\mathcal{F}}\) 的）称为 \(x\) 处截面 \(s\) 的芽。

设 \(a\) 为 B 的一点。集 \(\mathcal{B}(a)\) （B 的开集 \(\mathrm{U} \in \mathcal{B}\) （包含 \(a\) 者）所成、按关系 \(\supset\) 排序者）是过滤的。把 \(\mathcal{F}\) 的指标集限制到 \(\mathcal{B}(a)\)，我们得到一个直接系统 \(((\mathcal{F}(\mathrm{U}))_{\mathrm{U} \in \mathcal{B}(a)}, (r_{\mathrm{UV}}))\)。按定义 (E, III, p. 60)，这个系统的归纳极限是对 \((\mathrm{U}, s)\) （满足 \(a \in \mathrm{U}\) 与 \(s \in \mathcal{F}(\mathrm{U})\) 者）所成的集关于等价关系 \(R\) 的商，后者由 \(R((\mathrm{U},s),(\mathrm{U}',s'))\) 定义，当且仅当存在 \(\mathrm{W} \in \mathcal{B}\) ，它包含 \(a\) 、含于 \(\mathrm{U} \cap \mathrm{U}'\) 且使得 \(r_{\mathrm{WU}}(s) = r_{\mathrm{WU}'}(s')\)。因此，按归纳极限的定义，这个极限等同于茎 \(\mathcal{F}_a\) （\(\mathcal{F}\) 在 \(a\) 处的）。

设 \(\mathcal{G}\) 为 B 上的、相对于基 \(\mathcal{B}\) 的一个预层，并设 \(\varphi = (\varphi_{\mathrm{U}})_{\mathrm{U} \in \mathcal{B}}\) 为从 \(\mathcal{F}\) 到 \(\mathcal{G}\) 中的预层态射。映射 \((\mathrm{U}, s, x) \mapsto (\mathrm{U}, \varphi_{\mathrm{U}}(s), x)\) （从 \(\mathrm{X}_{\mathcal{F}}\) 到 \(\mathrm{X}_{\mathcal{G}}\) 中的）与等价关系 \(\mathrm{R}_{\mathcal{F}}\) 和 \(\mathrm{R}_{\mathcal{G}}\) 相容（由预层态射的定义）。用 \(\mathrm{E}(\varphi): \mathrm{E}_{\mathcal{F}} \to \mathrm{E}_{\mathcal{G}}\) 表示由它过渡到商而得的映射。对每个 \(\mathrm{U} \in \mathcal{B}\) 与每个 \(s \in \mathcal{F}(\mathrm{U})\) ，我们有

\[
\mathrm{E} (\varphi) \circ \sigma_ {\mathcal {F}} (\mathrm{U}, s) = \sigma_ {\mathcal {G}} (\mathrm{U}, \varphi_ {\mathrm{U}} (s));
\]

因此映射 \(\mathrm{E}(\varphi)\) 是连续的。映射 \(\mathrm{E}(\varphi)\) 是一个 B-态射；称之为从 \(E_{F}\) 到 \(E_{G}\) 中的、与预层态射 \(\varphi\) 相伴的 B-态射。对每个 \(a \in B\) ，\(\mathrm{E}(\varphi)\) 限制在 a 上的纤维上，给出从 F 的茎 \(F_{a}\) 到 G 的茎 \(G_{a}\) 中的一个映射；把它记作 \(\varphi_{a}\) 。它也是诸映射 \(\varphi_{U}\) 的归纳极限 (E, III, p. 63)，其中 U 取遍属于基 B 且包含 a 的开集所成的集 \(\mathcal{B}(a)\) 。

我们有 \(\operatorname{E}(\operatorname{Id}_{\mathcal{F}}) = \operatorname{Id}_{\operatorname{E}_{\mathcal{F}}}.\)

设 \(\mathcal{H}\) 为 B 上的、相对于 \(\mathcal{B}\) 的一个预层，并设 \(\psi = (\psi_{\mathrm{U}})\) 为从 \(\mathcal{G}\) 到 \(\mathcal{H}\) 中的预层态射。对 \([\mathrm{U}, s, x] \in \mathrm{E}_{\mathcal{F}}\) ，我们有

\[
\begin{array}{r l} & {\mathrm{E} (\psi \circ \varphi) ([ \mathrm{U}, s, x ]) = [ \mathrm{U}, \psi_ {\mathrm{U}} \circ \varphi_ {\mathrm{U}} (s), x ]} \\ & {\qquad = \mathrm{E} (\psi) ([ \mathrm{U}, \varphi_ {\mathrm{U}} (s), x ])} \\ & {\qquad = \mathrm{E} (\psi) \circ \mathrm{E} (\varphi) ([ \mathrm{U}, s, x ]).} \end{array}
\]

因此有 \(\mathrm{E}(\psi \circ \varphi) = \mathrm{E}(\psi) \circ \mathrm{E}(\varphi)\) 。特别地，若 \(a\) 是 B 的一点，则 \((\psi \circ \varphi)_a = \psi_a \circ \varphi_a\) 。

若 \(\varphi\) 是一个同构，则对 \(\mathrm{E}(\varphi)\) 同样成立。

注. --- 设 \(\mathcal{F}\) 为 B 上的、相对于基 \(\mathcal{B}\) 的一个层。设 B' 为 B 的一个开子集，设 \(\mathcal{B}'\) 为 B' 的拓扑的一个基，使得 \(\mathcal{B}' \subset \mathcal{B}\) 。设 \(F \mid B'\) 为 B' 上的、相对于基 \(B'\) 的、由 F 经限制而得的预层。

\begin{enumerate}
\def\labelenumi{\arabic{enumi})}
\tightlist
\item
  集合 \(X_{F|B'}\) 于是是 \(X_{F}\) 的一个子集，且 \(R_{F}\) 在 \(X_{F|B'}\) 上诱导的等价关系就是等价关系 \(R_{F|B'}\)。由此我们得到从 \(E_{F|B'}\) 到 \(E_{F}\) 中的一个典范单射 i。它的像是 \(p^{-1}(B')\)；对元素 \([U, s, x]\) （\(p^{-1}(B')\) 的每个），存在 \(B'\) 的一个元素 V，使得 \(x \in V\)、\(V \subset U\)，且有 \([U, s, x] = i([V, r_{\mathrm{VU}}(s), x])\) 。映射 i 是连续的，因为 \(X_{F|B'}\) 的拓扑是使得由 \(x \mapsto [U, s, x]\) （对 \(U \in B'\) 与 \(s \in \mathcal{F}(U)\) ）定义的映射都连续的最细拓扑。根据 I, p. 30 的命题 6 的推论 2，从 \(E_{F|B'}\) 到 \(E_{F}\) 中的典范单射 i 诱导从 \(E_{F|B'}\) 到 \(p^{-1}(B')\) 上的一个 B'-同构。
\end{enumerate}

特别地，当 B' 等于 B 时，\(i: E_{F|B'} \to E_F\) 是平展空间的一个 B-同构。

\begin{enumerate}
\def\labelenumi{\arabic{enumi})}
\setcounter{enumi}{1}
\tightlist
\item
  设 \(\mathcal{G}\) 为 B 上的、相对于基 \(\mathcal{B}\) 的一个预层，并设 \(\varphi\colon\mathcal{F}\to\mathcal{G}\) 为一个预层态射。图 \(\varphi'=(\varphi_{\mathrm{U}})_{\mathrm{U}\in\mathcal{B}'}\) 是从 \(\mathcal{F}\mid\mathcal{B}'\) 到 \(\mathcal{G}\mid\mathcal{B}'\) 中的预层态射。图
\end{enumerate}

\[
\begin{array}{c} \mathrm{E} _ {\mathcal {F} | \mathcal {B} ^ {\prime}} \xrightarrow {\mathrm{E} (\varphi^ {\prime})} \mathrm{E} _ {\mathcal {G} | \mathcal {B} ^ {\prime}} \\ \Biggl \downarrow_ {i} \\ \mathrm{E} _ {\mathcal {F}} \xrightarrow {\mathrm{E} (\varphi)} \mathrm{E} _ {\mathcal {G}}  , \end{array}
\]

其中 i 与 \(i'\) 是典范单射，是交换的。

例. --- 1) 设 B 为一个拓扑空间，\(\mathcal{B}\) 为 B 的拓扑的一个基，F 为一个集合。对 \(\mathcal{F}\) 取 B 上的、相对于 \(\mathcal{B}\) 的、由 \(\mathcal{F}(\mathrm{U}) = \mathrm{F}\) （对一切 \(\mathrm{U} \in \mathcal{B}\)）与 \(r_{\mathrm{UV}} = \mathrm{Id}_{\mathrm{F}}\) （对 \(\mathcal{B}\) 中满足 \(\mathrm{U} \subset \mathrm{V}\) 的每对元素 (U, V)）定义的预层。映射 \([\mathrm{U}, s, x] \mapsto (x, s(x))\) 是从 B-空间 \(\mathrm{E}_{\mathcal{F}}\) 到 B-空间 \(\mathrm{B} \times \mathrm{F}\) （其中 F 赋有离散拓扑）上的一个 B-同构。注意，当 \(\mathcal{B}\) 是 B 的开子集全体时，B 上的预层 \(\mathcal{F}\) 是层，仅当集 F 化为单点集时（参看 I, p. 43, 注）。

\begin{enumerate}
\def\labelenumi{\arabic{enumi})}
\setcounter{enumi}{1}
\item
  设 B 为一个拓扑空间，\((\mathrm{E},p)\) 为一个 B-空间。对 F 取 \((\mathrm{E},p)\) 的连续截面在 B 上的层。映射 \((\mathrm{U},s,x)\mapsto s(x)\) （从 \(X_{F}\) 到 E 中的）与等价关系 \(R_{F}\) 相容。由过渡到商得到的映射 e: \(E_{F}\to E\) 是一个 B-态射；称 B-态射 e 为典范的。e 的像是 p 在 B 的诸开子集上的连续截面的像的并。因此，若 p 是平展的，则映射 e 是满射的 (I, p. 33, prop. 9)。另一方面，映射 e 是单射的，当且仅当对 B 的每个开子集 U 和 p 在 U 上的每对连续截面 \((s,s')\) ，由点 \(x\in U\) （满足 \(s(x)=s'(x)\) 者）所成的集是开的；若 p 是平展的，则特别属于这种情形 (I, p. 34, prop. 11, b))。因此，若 \((\mathrm{E},p)\) 是一个平展 B-空间，则映射 e 是一个 B-同构。
\item
  设 B 为一个拓扑空间，\(\mathcal{B}\) 为 B 的拓扑的一个基，\(\mathcal{F}\) 为 B 上的、相对于 \(\mathcal{B}\) 的一个预层，\(\mathcal{L}\) 为 \(\mathcal{F}\) 的一个子预层。那么，集 X \(_{\mathcal{L}}\) 含于集 X \(_{\mathcal{F}}\) ，且等价关系 R \(_{\mathcal{F}}\) 在 X \(_{\mathcal{L}}\) 上诱导等价关系 R \(_{\mathcal{L}}\) 。因此，B-态射 E(i): E \(_{\mathcal{L}} \to\) E \(_{\mathcal{F}}\) （与典范态射 i: \(\mathcal{L} \to \mathcal{F}\) (I, p. 48, 例 2) 相伴者）是单射的。由于 E \(_{\mathcal{L}}\) 与 E \(_{\mathcal{F}}\) 是平展 B-空间，映射 E(i) 是开的甚至是平展的 (I, p. 30, cor. 1)，从而诱导从 E \(_{\mathcal{L}}\) 到 E \(_{\mathcal{F}}\) 的一个开子集上的一个同胚。
\end{enumerate}

